\chapter{Introduction}\label{ch:introduction}
% Background, problem statement, purpose and scope of the project.

% ------------------------------------------------------------
% EXAMPLE of inserting a DrawIO diagram:
% 1. Save/edit the diagram in  assets/drawio/<name>.drawio
% 2. Reference it with \drawiofig - the PDF is generated
%    AUTOMATICALLY by latexmk when you compile
%    (requires draw.io Desktop locally).
%
% Usage:          \drawiofig{file-name}{Caption}{fig:label}
% Optional width: \drawiofig[0.5\linewidth]{...}{...}{...}
% NOTE: all three arguments are required!
% ------------------------------------------------------------
\drawiofig{v-model}{Example: DrawIO diagram inserted automatically (the V-model)}{fig:vmodel-example}

As shown in \autoref{fig:vmodel-example}, a diagram is inserted simply by
referencing the name of the .drawio file.


\section{Conceptual overview}\label{sec:conceptual-overview}
\autoref{fig:rich-picture} shows who uses PrepMeUp and what moves between them.
The household talks to the system through the customer app, the administrator
through the admin interface, and the supplier and the payment are simulated
inside the project.

\drawiofig[\linewidth]{rich-picture}{Rich picture of PrepMeUp. The household only meets the system through the app and the delivered packages; the dashed boxes are simulated within the project, and the screens are illustrative.}{fig:rich-picture}
